Для реализации алгоритма 2-аппроксимации решения TSP с использованием минимального остовного дерева выбран язык программирования Rust, так как он предоставляет высокою степень оптимизации, является компилируемым и безопасно работает с памятью.

\subsection{Построение MST. Алгоритм Прима}

\begin{figure}
	\centering
	\begin{tikzpicture}[
			node distance=2cm,
			block/.style={draw, rectangle, rounded corners, align=center, minimum width=3cm, minimum height=1cm},
			arrow/.style={->, thick}
		]

		\node[block] (start) {Начало};
		\node[block, below of=start] (init) {Выбрать \\ начальную вершину};
		\node[block, below of=init] (select) {Найти минимальное \\ ребро, соединяющее \\ дерево и новую вершину};
		\node[block, below of=select] (add) {Добавить ребро \\ в MST};
		\node[block, below of=add] (check) {Все вершины \\ добавлены?};
		\node[block, below of=check] (end) {Конец};

		\draw[arrow] (start) -- (init);
		\draw[arrow] (init) -- (select);
		\draw[arrow] (select) -- (add);
		\draw[arrow] (add) -- (check);
		\draw[arrow] (check) -- node[right]{Да} (end);
		\draw[arrow] (check.east) -- ++(2,0) |- node[right]{Нет} (select.east);

	\end{tikzpicture}
	\caption{Общая схема алгоритма Прима}
\end{figure}

Для построения минимального остовного дерева, программа использует алгоритм Прима, основанный на использовании минимальной кучи (min-heap). Этот алгоритм выбран так как он быстрее работает на плотных графах, по условию задачи, граф является полным --- а значит, максимально плотным.

Описаие алгоритма:
\begin{enumerate}
	\item создаем массивы \texttt{key}, \texttt{parent}, \texttt{in\_mst}, для хранения соответственно минимального веса ребра, соединяющего вершину с MST, индекса ``родительской'' вершины в MST и флага, определяющего включена ли вершина в MST. Также создадим минимальную кучу;
	\item установим для нулевой вершины \texttt{key} в $0$, а \texttt{parent} --- \texttt{None};
	\item добавим все вершины графа в кучу;
	\item пока в куче есть вершины:
	      \begin{itemize}
		      \item сохраним текущую вершину в переменной \texttt{u}
		      \item если вершина уже добавлена в MST, переходим к следующей вершине;
		      \item добавляем вершину в MST, установив \texttt{in\_mst} в \texttt{true};
		      \item для всех вершин \texttt{v} в графе:
		            \begin{itemize}
			            \item если \texttt{v} уже  добавлена в MST, переходим к следующей;
			            \item если ребро $(u, v)$ имеет меньший вес, чем \texttt{key[v]}, то установим \texttt{key[v]} равным весу этого ребра, \texttt{parent[v]} равным \texttt{u} и добавим вершину \texttt{v} в кучу с новым значением \texttt{key}.
		            \end{itemize}
	      \end{itemize}
\end{enumerate}

\begin{lstlisting}[caption=Реализация алгоритма Прима на Rust]
    let mut key = vec![f32::INFINITY; vertices.len()];
    let mut parent = vec![None; vertices.len()];
    let mut in_mst = vec![false; vertices.len()];

    let mut heap = MyPriorityQueue::<HeapItem>::new();

    // initialization
    key[0] = 0.0;
    parent[0] = None;

    for (key, index) in key.iter().zip(0..) {
        heap.push(HeapItem(*key, index));
    }

    while let Some(HeapItem(_key_u, u)) = heap.pop() {
        if in_mst[u] {
            continue;
        }
        in_mst[u] = true;

        for v in 0..vertices.len() {
            if in_mst[v] || v == 0 {
                continue;
            }

            let edge_length = vertices[u].distance(vertices[v]);
            if edge_length < key[v] {
                key[v] = edge_length;
                parent[v] = Some(u);

                heap.push(HeapItem(key[v], v));
            }
        }
    }

\end{lstlisting}

\begin{figure}[h]
	\centering
	\begin{tikzpicture}[scale=1.2]
		\tikzstyle{vertex}=[circle, draw, minimum size=7mm]
		\tikzstyle{edge}=[draw]
		\tikzstyle{mst}=[draw, thick, color=green]

		\node[vertex] (A) at (0,0) {A};
		\node[vertex] (B) at (3,1) {B};
		\node[vertex] (C) at (3,-1) {C};
		\node[vertex] (D) at (6,0) {D};
		\node[vertex] (E) at (6, -2) {E};
		\node[vertex] (F) at (9, -1) {F};
		\node[vertex] (G) at (9, 1) {G}

		\draw[edge] (A) -- (B);
		\draw[edge] (A) -- (C);
		\draw[edge] (B) -- (C);
		\draw[edge] (B) -- (D);
		\draw[edge] (C) -- (D);
		\draw[edge] (D) -- (G);
		\draw[edge] (C) -- (E);
		\draw[edge] (C) -- (F);
		\draw[edge] (E) -- (F);
		\draw[edge] (D) -- (F);

		\draw[mst] (A) -- (B);
		\draw[mst] (B) -- (D);
		\draw[mst] (B) -- (C);
		\draw[mst] (C) -- (E);
		\draw[mst] (D) -- (F);
		\draw[mst] (D) -- (G);

	\end{tikzpicture}
	\caption{Пример построения минимального остовного дерева}
\end{figure}

После завершения алгоритма, список смежности MST восстанавливается по информации из массива \texttt{parent}.

\subsection{Построение итогового обхода}

По списку смежности MST, производится обход в глубину, вершины записываются в порядке посещения.

Если вершина уже посещена ранее, ее можно пропустить, напрямую соединив текущую вершину со следующуй непосщенной. Благодаря неравенству треугольника длина маршрута при этом не увеличивается.

\begin{figure}[h]
	\centering
	\begin{tikzpicture}[
			level distance=1.6cm,
			every node/.style={circle, draw},
			level 1/.style={sibling distance=3cm},
			level 2/.style={sibling distance=3cm}
		]
		\node{A}
		child {
				node {B}
				child {
						node {C} child { node {E} }
					}
				child {
						node {D}
						child {
								node {G}
							}
						child {
								node {F}
							}
					}
			}
	\end{tikzpicture}
	\caption{Обход дерева в глубину}
\end{figure}

\begin{figure}[h]
	\centering
	\begin{tikzpicture}[
		every node/.style={circle, draw},
		node distance=1.5cm,
		tour/.style={
		-{Stealth[length=3mm, width=2mm]},
		thick,
		green,
		}
		]

		\node (a) {A};
		\node[right of=a] (b) {B};
		\node[right of=b] (c) {C};
		\node[right of=c] (e) {E};
		\node[right of=e] (c2) {C};
		\node[right of=c2] (b2) {B};
		\node[right of=b2] (d) {D};
		\node[right of=d] (g) {G};
		\node[right of=g] (d2) {D};
		\node[right of=d2] (f) {F};
		\node[right of=f] (a2) {A};

		\draw[->] (a) -- (b);
		\draw[->] (b) -- (c);
		\draw[->] (c) -- (e);
		\draw[->] (e) -- (c2);
		\draw[->] (c2) -- (b2);
		\draw[->] (b2) -- (d);
		\draw[->] (d) -- (g);
		\draw[->] (g) -- (d2);
		\draw[->] (d2) -- (f);

		\draw[tour] (a.north) .. controls +(up:7mm) and +(up:7mm) .. (b.north);
		\draw[tour] (b.south) .. controls +(down:7mm) and +(down:7mm) .. (c.south);
		\draw[tour] (c.north) .. controls +(up:7mm) and +(up:7mm) .. (e.north);
		\draw[tour] (e.south) .. controls +(down:7mm) and +(down:7mm) .. (d.south);
		\draw[tour] (d.north) .. controls +(up:7mm) and +(up:7mm) .. (g.north);
		\draw[tour] (g.south) .. controls +(down:7mm) and +(down:7mm) .. (f.south);
		\draw[tour] (f) -- (a2);

	\end{tikzpicture}
	\caption{Построение гамильтонова обхода на основе обхода в глубину}
\end{figure}

В итоге, получаем гамильтонов цикл, имеющий вес, максимум, в два раза больший веса решения TSP на заданном графе.
